\chapter{Herglotz arithmetic and quadratic regulators}
\label{ch:09-herglotz}

The Herglotz function joins the digamma series to finite dilogarithm sums.
At rational arguments this supplies exact evaluations; at quadratic units it
connects continued-fraction cycles with real-quadratic Kronecker limits.
The same experimental strategy thus leads from a finite cyclotomic sum to
regulator candidates in class fields.

\section{The function and the exact dilogarithm sum}
For $x>0$ define
\[
F(x)=\sum_{n\ge1}\frac{\psi(nx)-\log(nx)}n,
\qquad
J(x)=\int_0^1\frac{\log(1+t^x)}{1+t}\,dt.
\]
The series is locally uniformly convergent on the positive line and extends
analytically to the cut plane $\mathbb C\setminus(-\infty,0]$.
Radchenko--Zagier relate the two functions by

\begin{equation}\label{herg:eq:Jint}
  J(x)=\int_0^1\frac{\log(1+t^{x})}{1+t}\,dt,\qquad J(x)=F(2x)-2F(x)+F(x/2)+\frac{\pi^2}{12x},
\end{equation}
which relates the companion $J$ to the Herglotz function $F$. For the rational-argument results the engine is
Radchenko--Zagier's Theorem~3, an \emph{exact, finite} evaluation obtained from the second integral
representation of $F$ via the distribution relations:
\begin{equation}\label{herg:eq:dilogsum}
  F\!\Bigl(\tfrac PQ\Bigr)-F(1)=\sum_{\alpha^P=1}\Bigl(\sum_{\substack{\beta^Q=1\\\beta\neq1}}f(\alpha,\beta)
  -\sum_{\substack{\beta^P=1\\\beta\neq1}}f(\alpha,\beta)\Bigr),\quad
  f(\alpha,\beta)=\Li_2\!\Bigl(\tfrac{\beta}{\beta-1}\Bigr)-\Li_2\!\Bigl(\tfrac{\alpha-\beta}{1-\beta}\Bigr).
\end{equation}
The finite sum is an exact analytic identity, with compatible principal
branches at its arguments. Floating-point evaluation still has rounding
and cancellation errors. It is a useful independent comparator for the
defining integral, not a zero-error numerical oracle.

\section{Quadratic-unit values: Tables 1 \& 2, and two errata}\label{herg:sec:tables}
For real quadratic units the relevant object is the companion
$\Jcal(x)=J(x)-\tfrac{\log^2 2}{2}+\tfrac{\pi^2}{24}(x-\tfrac1x)$, and the arithmetic atoms are
$S_{p,q}=\log\eps_p\,\log\eps_q$ where $\eps_d$ is the fundamental unit of $\mathbb Q(\sqrt d)$ for a
fundamental discriminant $d>1$ and $\eps_1:=2$. Tables~1 and~2 of the source paper express, for the $45$
known one-class-per-genus cases, the quantities
$4\bigl(\Jcal(n+\sqrt{n^2-1})-n\tfrac{\pi^2}{24}\bigr)$ and
$2\bigl(\Jcal(n+\sqrt{n^2+1})-n\tfrac{\pi^2}{24}\bigr)$ as integer combinations of the $S_{p,q}$.

We re-implemented $\Jcal$ and $S_{p,q}$ in PARI/GP (\texttt{intnum} integrals, \texttt{bnfinit} regulators)
and evaluated both sides of all $45$ rows at $200$-digit precision. All $16$ rows of Table~2 and $27$ of the
$29$ rows of Table~1 agree to better than $10^{-80}$. The two failures are $O(1)$:
\[
  \text{$n=43$:}\ \ \text{LHS}-\text{RHS}=-13.0160\ldots,\qquad
  \text{$n=61$:}\ \ \text{LHS}-\text{RHS}=-2.9011\ldots .
\]
Each resolves to a single exact correction, confirmed against the high-DPI scan to be a typo in the printed
paper rather than an OCR artifact.

\begin{proposition}[Two corrections to Table 1]
\leavevmode
\begin{itemize}[topsep=2pt,itemsep=1pt]
\item $n=43$ ($n^2-1=2^3\cdot3\cdot7\cdot11$): the printed right-hand side overshoots by exactly $S_{33,56}$
$\bigl((\text{LHS}-\text{RHS})/S_{33,56}=-1.000\ldots$ to $110$ digits$\bigr)$; the spurious $+S_{33,56}$
term should be deleted.
\item $n=61$ ($n^2-1=2^3\cdot3\cdot5\cdot31$): the printed $S_{60,1032}$ should read $S_{60,248}$. The
printed subscript $1032=2^3\cdot3\cdot43$ carries the prime $43$, foreign to $n^2-1=3720$, whereas
$60\cdot248=14880=4\cdot3720$ is the correct genus splitting.
\end{itemize}
A brute-force search over all single edits (delete a term, change a coefficient in $[-4,4]$, or replace a
subscript by any admissible fundamental-discriminant pair) returns these as the \emph{unique} fixes. With
both applied, all $45$ rows verify to better than $10^{-80}$.
\end{proposition}

The quadratic-unit tables concern specified arithmetic families. They do
not imply that every even-trace unit has a pure-logarithm evaluation:
outside genus characters the Kronecker-limit formula introduces additional
class-group coordinates. Those coordinates are developed below.

\section{Rational arguments: a family of theorems}\label{herg:sec:rational}
The paper writes out only two rational-argument cases: $F(n)$ (pure logarithms, its eq.~23) and the lone
fractional example $F(2/5)$ (its eq.~24). We now give the general laws.

\subsection{Two reductions on the line \texorpdfstring{$\Real=\tfrac12$}{Re=1/2}}
Everything rests on one observation: for a unit $\beta=e^{i\theta}$,
\begin{equation}\label{herg:eq:line}
  \frac{\beta}{\beta-1}\ \text{has real part}\ \tfrac12,\qquad
  \arg\frac{\beta}{\beta-1}=\frac{\theta}{2}-\frac{\pi}{2},\qquad
  \Bigl|\frac{\beta}{\beta-1}\Bigr|=\frac{1}{2\sin(\theta/2)} .
\end{equation}
By the reflection $\Li_2(z)+\Li_2(1-z)=\tfrac{\pi^2}{6}-\log z\log(1-z)$ with $1-z=\bar z$ on $\Real z=\tfrac12$,
\begin{equation}\label{herg:eq:reRe}
  2\Real\Li_2\Bigl(\tfrac12+iy\Bigr)=\frac{\pi^2}{6}-\log^2|z|-(\arg z)^2 .
\end{equation}
The second reduction handles the $\alpha=-1$ part of~\eqref{herg:eq:dilogsum}: $\tfrac{-1-\beta}{1-\beta}=-i\cot\tfrac\theta2$
is purely imaginary, and for real $y$
\begin{equation}\label{herg:eq:imax}
  \Real\Li_2(iy)=\tfrac14\Li_2(-y^2).
\end{equation}

\subsection{$F(1/q)$: pure logarithms}
\begin{theorem}\label{herg:thm:F1q}
For all $q\ge2$, with $m=\lfloor(q-1)/2\rfloor$,
\[
  F\!\Bigl(\tfrac1q\Bigr)-F(1)=-\frac{(q-1)(3q-2)}{24q}\,\pi^2
  -\sum_{k=1}^{m}\log^2\!\Bigl(2\sin\tfrac{\pi k}{q}\Bigr)\;-\;[\,q\ \mathrm{even}\,]\,\tfrac12\log^2 2 .
\]
\end{theorem}
\begin{proof}
With $P=1$ only $\alpha=1$ occurs and the inner $\beta^P$-sum is empty, so~\eqref{herg:eq:dilogsum} reads
$F(1/q)-F(1)=\sum_{\beta^q=1,\beta\neq1}(\Li_2(\tfrac{\beta}{\beta-1})-\tfrac{\pi^2}{6})$. Pairing $\beta$ with
$\bar\beta$ and applying~\eqref{herg:eq:reRe} gives the $\pi^2$ coefficient
$-\bigl[\tfrac{q-1}{12}+\sum_{k=1}^m(\tfrac kq-\tfrac12)^2\bigr]=-\tfrac{(q-1)(3q-2)}{24q}$ and the log-sine sum.
For even $q$ the self-conjugate root $\beta=-1$ contributes $\Li_2(\tfrac12)=\tfrac{\pi^2}{12}-\tfrac12\log^2 2$.
\end{proof}

\subsection{$F(2/q)$: the general theorem}
\begin{theorem}\label{herg:thm:F2q}
For every $q\ge3$,
\[
  F\!\Bigl(\tfrac2q\Bigr)-F(1)=-\frac{(q-1)(q-2)}{12q}\pi^2+\log^2 2
  -2\!\!\sum_{k=1}^{\lfloor(q-1)/2\rfloor}\!\!\log^2\!\Bigl(2\sin\tfrac{\pi k}{q}\Bigr)
  -\tfrac12\!\!\sum_{k=1}^{\lfloor(q-1)/2\rfloor}\!\!\Li_2\!\Bigl(-\cot^2\tfrac{\pi k}{q}\Bigr),
\]
where for even $q$ the two sums run to $k=\tfrac q2-1$ and a further $-\log^2 2$ is added.
\end{theorem}
\begin{proof}[Proof (odd $q$)]
Here $\alpha\in\{1,-1\}$ and the inner $\beta^P$-sum is the single term $\beta=-1$. For the first part,
$f(1,\beta)+f(-1,\beta)=2\Li_2(\tfrac{\beta}{\beta-1})-\tfrac{\pi^2}{6}-\Li_2(\tfrac{-1-\beta}{1-\beta})$.
Summing over $\beta\neq1$ and using~\eqref{herg:eq:reRe} for the first pieces and~\eqref{herg:eq:imax}
(with $\tfrac{-1-\beta}{1-\beta}=-i\cot\tfrac{\pi k}q$, pairing $k\leftrightarrow q-k$) for the last gives
$2S_1-(q-1)\tfrac{\pi^2}{6}-\tfrac12\sum_{k}\Li_2(-\cot^2\tfrac{\pi k}q)$, where
$S_1=\sum_{k=1}^{(q-1)/2}[\tfrac{\pi^2}{6}-\log^2(2\sin\tfrac{\pi k}q)-\pi^2(\tfrac kq-\tfrac12)^2]$. The boundary
gives $-\sum_\alpha f(\alpha,-1)=+\log^2 2$ from $\Li_2(\tfrac12)$. Collecting, the $\tfrac{\pi^2}{6}$ terms
cancel and the $\pi^2$ coefficient is $-2\sum_{k=1}^{(q-1)/2}(\tfrac kq-\tfrac12)^2$, which equals
$-\tfrac{(q-1)(q-2)}{12q}$ by an elementary identity (verified symbolically). The even-$q$ case differs only
in that $\beta=-1$ is now a genuine $q$-th root, contributing the extra $-\log^2 2$.
\end{proof}
The algebraic dilogarithm arguments make the field structure explicit.
At $q=5$ the symmetric sum can be reduced further as described below.
Unsuccessful searches at larger conductors do not prove that no other
closed form exists; the finite-sum theorem remains valid at every $q$.

\subsection{The $F(n/(n+1))$ family and the $F(p/2)$ reciprocity}
\begin{theorem}\label{herg:thm:np1}
For all $n\ge2$, with $S(q)=\sum_{k=1}^{q-1}\log^2(2\sin\tfrac{\pi k}q)$,
\[
  F\!\Bigl(\tfrac{n}{n+1}\Bigr)-F(1)=-\frac{3n^2+3n+2}{24n(n+1)}\pi^2+\Li_2\!\Bigl(\tfrac{n}{n+1}\Bigr)
  +\tfrac12\log^2\!\tfrac{n+1}{n}-\tfrac12 S(n+1)+\tfrac12 S(n),
\]
the single dilogarithm $\Li_2(n/(n+1))$ appearing with coefficient exactly $+1$ (the source paper's
``bilinear-log'' Example~2, made explicit). The $\pi^2$ coefficient equals $-\tfrac18-\tfrac1{12n(n+1)}$.
\end{theorem}
\begin{theorem}[Reciprocity]\label{herg:thm:p2}
For all odd $p\ge3$,
\[
  F\!\Bigl(\tfrac p2\Bigr)+F\!\Bigl(\tfrac2p\Bigr)-2F(1)=-\frac{(p-2)^2}{12p}\pi^2+\tfrac12\log^2\!\tfrac p2 .
\]
The $\Li_2(-\cot^2\frac{\pi k}{p})$ terms of $F(p/2)$ are exactly minus those of $F(2/p)$ and cancel in the
sum; combined with Theorem~\ref{herg:thm:F2q} this yields $F(p/2)$ in closed form.
\end{theorem}

\subsection{The derivative law}
\begin{theorem}[Radchenko--Zagier Proposition 4, $P=2$]\label{herg:thm:deriv}
For all $q\ge3$, with $T(q):=\tfrac2q F'(2/q)-(1+\log2)$,
\[
  T(q)=\frac{q^2-4}{24q}\pi^2+W(q),\qquad
  W(q)=\sum_{k=1}^{q-1}C(2k/q)\,\Cl_2\!\Bigl(\tfrac{2\pi k}q\Bigr),
\]
where $C(t)=\cot(\pi t)$ for $t\notin\mathbb Z$ and $C(t)=0$ for
$t\in\mathbb Z$, with an extra $-\log2$ for even $q$.
This assigned value at integers is part of the formula; a literal cotangent
there would be undefined. The sum provides a Clausen representation,
without an independence assertion.
\end{theorem}

\subsection{Rational-argument samples and the companion integral}
The source's proposed criterion $\operatorname{rad}(q)\mid30$ is false
as a necessary condition for a reduction to rational dilogarithms:
Theorem~\ref{herg:thm:F1q} gives a pure-logarithm answer for $F(1/7)$.
Numerator, denominator and the available functional equations must all be
considered. No replacement universal classification is asserted here.
The following experimentally recovered samples are useful:
\begin{align*}
F(3/4)-F(1)&=-\frac{19}{144}\pi^2-\frac34\log^22
-2\log2\log3+\frac74\log^23+2\Li_2(1/3),\\
F(4/5)-F(1)&=\frac3{80}\pi^2+\frac{11}4\log^22-\frac12\log^25
-\Li_2(1/5)-\log^2(2\sin(\pi/5))-\log^2(2\sin(2\pi/5)),\\
J(2/3)&=-\frac{\pi^2}{144}+\frac34\log^22,\\
J(2/5)&=\frac{11\pi^2}{240}+\frac34\log^22-2\log^2\phi.
\end{align*}
The last two values follow by combining the $F$ evaluations through
\eqref{herg:eq:Jint}. Their existence does not prove that these are the
only $q$ for which $J(2/q)$ simplifies.

\section{Bridges to the golden ladders and the Clausen lattice}\label{herg:sec:bridges}
The new constants are not isolated; they touch two corpora we have already mapped.

\paragraph{The golden bridge ($q=5$).} The arguments $-\cot^2\tfrac\pi5=-\tfrac{5+2\sqrt5}5$ and
$-\cot^2\tfrac{2\pi}5=-\tfrac{5-2\sqrt5}5$ lie in $\mathbb Q(\sqrt5)$, the field of the golden polylogarithm
ladders. Comparing Theorem~\ref{herg:thm:F2q} ($q=5$) with eq.~24 gives the functional equation
\[
  \Li_2\!\Bigl(\tfrac45\Bigr)+\Li_2\!\Bigl(-\cot^2\tfrac{\pi}{5}\Bigr)+\Li_2\!\Bigl(-\cot^2\tfrac{2\pi}{5}\Bigr)
  =2\log^2 2-2\log2\log\tfrac25-6\bigl(\log^2(2\sin\tfrac\pi5)+\log^2(2\sin\tfrac{2\pi}5)\bigr),
\]
hence the golden-sum closed form (with $\varphi=\tfrac{1+\sqrt5}2$)
\[
  \Li_2\!\Bigl(-\cot^2\tfrac{\pi}{5}\Bigr)+\Li_2\!\Bigl(-\cot^2\tfrac{2\pi}{5}\Bigr)
  =\Li_2\!\Bigl(\tfrac15\Bigr)-\frac{\pi^2}{6}-3\log^2\varphi+\tfrac14\log^2 5 .
\]
The symmetric sum has the displayed reduction. The antisymmetric
combination $X=\frac12[\Li_2(-\cot^2(\pi/5))-
\Li_2(-\cot^2(2\pi/5))]$ did not reduce in the source's tested basket.
Its independence is open. A formal symbol with nonzero Bloch boundary
is not subject to a vanishing-rank conclusion for the Bloch group.

\paragraph{The Clausen bridge.}
The coefficients of $W(q)$ are real algebraic numbers contained in
$\mathbb Q(\zeta_{\operatorname{lcm}(4,q)})^+$; at odd conductors the factor
$i$ implicit in the cotangent can require an extension of
$\mathbb Q(\zeta_q)^+$. This is an algebraic coefficient issue, not a proof
of numerical period independence. Applying the distribution identities
of Chapter~\ref{ch:02-cyclotomic} gives
\[
W(8)=2\Cl_2(\pi/4)-2\Cl_2(3\pi/4)=G,
\qquad
W(12)=\frac{11\sqrt3}{9}\Cl_2(\pi/3).
\]
For the second formula, pair $k$ with $12-k$:
\[
W(12)=2\sqrt3\left[\Cl_2(\pi/6)-\Cl_2(5\pi/6)\right]
+\frac2{\sqrt3}\left[\Cl_2(\pi/3)-\Cl_2(2\pi/3)\right].
\]
Use $\Cl_2(\pi/6)-\Cl_2(5\pi/6)=V/2$ and
$\Cl_2(2\pi/3)=2V/3$. This proves the original report's coefficient
$22\sqrt3/9$ multiplying the first difference. It is also checked against
direct derivative quadrature in the fresh verification battery.

\section{Quadratic units and Kronecker-limit coordinates}
Define the normalized functions
\[
\mathcal F(x)=F(x)-F(1)+\frac{\pi^2}{12}(x-2+x^{-1})-\frac14\log^2x,
\quad
\mathcal J(x)=J(x)-\frac12\log^22+\frac{\pi^2}{24}(x-x^{-1}).
\]
For a real quadratic order of discriminant $D$, a narrow ideal class
$\mathcal B$ has a partial zeta function. Its Kronecker constant is
\[
\varrho(\mathcal B)=\lim_{s\to1}
\left[D^{s/2}\zeta(\mathcal B,s)-\frac{\log\varepsilon_D}{s-1}\right],
\]
where $\varepsilon_D>1$ is the least norm-one unit. In the reduced-form cycle
representation, $w_j>1>w_j'>0$, Radchenko--Zagier give
\begin{equation}\label{stark:eq:rho}
\varrho(\mathcal B)=2\gamma\log\varepsilon_D
-\ell(\mathcal B)\frac{\pi^2}6+
\sum_{j\bmod\ell}\left[\mathcal F(w_j)-\mathcal F(w_j')
+L(w_j'/w_j)\right],
\end{equation}
where $L$ is the unnormalized Rogers dilogarithm.
Class characters turn the $\varrho$ coordinates into Hecke $L$-values.
Their Theorem~5 expresses $\mathcal J$ at specified even-trace units in
these coordinates. The source's independent cycle engine reported checks
at $n=2,4,6,8$ with residuals about $5\times10^{-38}$; the engine was not
imported into this checkout.

\section{Regulator candidates and field degrees}
The published Stark examples and the source's follow-up fits are
conjectural or numerical evaluations. A cubic polynomial over $\mathbb Q$
defines a degree-three field $E$. It cannot itself be a Hilbert or ring
class field containing a quadratic field $K$. In the examples at issue,
its degree-six splitting field can be a cyclic cubic extension of $K$;
$E$ is a cubic subfield. The distinction corrects a repeated field-degree
error in the report.

For $u=16+\sqrt{257}$ the source reports
\begin{equation}\label{stark:eq:257}
\mathcal J(u)\overset{\rm num}=4\zeta(2)+\log2\log u+2\reg(E),
\quad E=\mathbb Q[x]/(x^3-2x^2-3x+1).
\end{equation}
The determinant of the two selected unit-logarithm columns is signed;
with the source's ordering it is $-\reg(E)$.
Vertical bars around a matrix denote its determinant, not an absolute
value. Replacing this negative determinant by its absolute value without
changing its external coefficient reverses the regulator term.

The other published example concerns $u=10+3\sqrt{11}$ and the order of
discriminant $396$. Its unit-logarithm expression is a Stark minor rather
than a full field regulator. We retain that structural distinction; the
unimported algebraic-unit definitions do not provide a locally replayable
formula for the minor.

\begin{table}[ht]\centering\small
\begin{tabular}{rrlccc}
\toprule
$n$&$D_0$&cubic field polynomial&$c_1$&$c_2$&$c_3$\\
\midrule
6&37&$x^3-x^2-3x+1$&$3/2$&$2/3$&$2/3$\\
10&101&$x^3-x^2-5x-1$&$5/2$&$2/3$&$2/3$\\
14&197&$x^3-x^2-7x-3$&$7/2$&$2/3$&$2/3$\\
16&257&$x^3-x^2-4x+3$&$4$&$1$&$2$\\
26&677&$x^3-x^2-11x-7$&$13/2$&$2/3$&$2/3$\\
\bottomrule
\end{tabular}
\caption{Historical numerical candidates
$\mathcal J(n+\sqrt{n^2+1})=c_1\zeta(2)+c_2\log2\log u+c_3\reg(E)$.
The source reports residuals below $1.2\times10^{-109}$. These are not
unconditional theorems.}\label{stark:tab:family}
\end{table}
The source reports that $n=6,10,14,16,26$ are the narrow-class-number-three
cases in its scan of $2\le n\le5000$. This is a bounded census, not a proof
that the family is finite. Negative fits at class numbers five and seven
likewise do not establish a universal ``single regulator if and only if
cyclic cubic'' theorem.

There is a precise conditional route to a leading-$L$ interpretation.
Suppose a non-Galois totally real cubic field $E$ has quadratic resolvent
$K$, and its normal closure supplies a nontrivial order-three class
character $\chi$ of $K$. Artin induction gives
$\zeta_E(s)=\zeta(s)L_K(s,\chi)$.
The analytic class-number formula then gives
\[
L_K(s,\chi)=h_E\reg(E)s^2+O(s^3),
\]
because $\zeta_E(s)=-h_E\reg(E)s^2/2+O(s^3)$ and $\zeta(0)=-1/2$.
Thus the coefficient equals $\reg(E)$ only after $h_E=1$ and the
field-character identification have been established. This argument does
not convert the numerical Herglotz fits into proofs of Stark's conjecture.
